\documentclass[11pt]{article}
\usepackage{amsmath,amssymb,amsthm,mathtools}
\newtheorem{theorem}{Theorem}
\newtheorem{lemma}{Lemma}
\newcommand{\E}{\mathbb E}
\newcommand{\Inf}{\operatorname{Inf}}
\newcommand{\cH}{\mathsf H}
\newcommand{\cU}{\mathsf U}
\title{A retained-difference entropy telescope for rare restriction trees}
\author{B05 / A023}
\date{13 October 2026}
\begin{document}
\maketitle

All logarithms are base two unless explicitly written as $\ln$.  For a
real function $r$ of squared norm at most one, put
\[
  \mathcal L(r)=\sum_S \widehat r(S)^2
       \log_2\frac1{\widehat r(S)^2},
\]
with $0\log(1/0)=0$.  For Boolean $f$, write $\cH(f)=\mathcal L(f)$,
$\pi_f(S)=\sum_{T\supseteq S}\widehat f(T)^2$, and
\[
  \cU(f)=\sum_S\widehat f(S)^2
       \log_2\frac{\pi_f(S)}{\widehat f(S)^2},\qquad
  V(f)=\cH(f)-\cU(f)\geq0.
\]

\section{One-node inequality}

Let $g,h:\{-1,1\}^n\to\{-1,1\}$ and let the fresh physical selector
$z$ choose between them.  Write
\[
 F(z,x)=\begin{cases}g(x),&z=-1,\\h(x),&z=1,\end{cases}
 \qquad d=\frac{g-h}{2},\qquad q=\E d^2=\Pr[g\ne h].
\]

\begin{theorem}[retained-difference node bound]\label{thm:node}
For every such genuine Boolean pair, including zero Fourier atoms,
\[
 \cU(F)-\frac{\cU(g)+\cU(h)}2
 \leq -V(F)+\frac{V(g)+V(h)}2
       +\mathcal L(d)+\frac{3}{\ln2}\,q.                 \tag{1}
\]
The coefficient $3/\ln2<4.329$ is absolute.
\end{theorem}

\begin{proof}
Put $a_S=\widehat g(S)$, $b_S=\widehat h(S)$,
$u_S=(a_S+b_S)/2$, and $v_S=(a_S-b_S)/2=\widehat d(S)$.
Then $M_S=u_S^2+v_S^2=(a_S^2+b_S^2)/2$ is a probability law.  Directly
splitting the Fourier atoms according to whether they contain $z$ gives
the exact identity
\[
 \cH(F)-\frac{\cH(g)+\cH(h)}2
  =\operatorname{JS}(a^2,b^2)
   +\sum_S M_S h_2\!\left(\frac{v_S^2}{M_S}\right),       \tag{2}
\]
where a zero summand is omitted.

We use the scalar inequality
\[
 \frac12\left[x^2\log_2\frac{2x^2}{x^2+y^2}
 +y^2\log_2\frac{2y^2}{x^2+y^2}\right]
 \leq\frac{(x-y)^2}{2\ln2}.                              \tag{3}
\]
For completeness, signs can only enlarge the right side, so take $x,y\geq0$;
by homogeneity set $y=1$ and $r=x$.  After multiplying by $2\ln2$, the
difference between the right and left sides is
\[
 G(r)=(r-1)^2-r^2\ln\frac{2r^2}{r^2+1}
                 -\ln\frac2{r^2+1}.
\]
Here
\[
 G'(r)=2(r-1)-2r\ln\frac{2r^2}{r^2+1}.
\]
The function
$1-r^{-1}-\ln(2r^2/(r^2+1))$ has derivative
$(r-1)^2/[r^2(r^2+1)]$ and vanishes at $r=1$.  Hence $G'$ has the sign of
$r-1$, and $G(r)\geq G(1)=0$.  Summing (3) yields
\[
 \operatorname{JS}(a^2,b^2)
 \leq\frac{\lVert a-b\rVert_2^2}{2\ln2}
 =\frac{2}{\ln2}q.                                      \tag{4}
\]

For the second term in (2), $M_S\leq1$ and
\[
 \begin{aligned}
 M_Sh_2(v_S^2/M_S)
 &=v_S^2\log_2\frac{M_S}{v_S^2}
   +u_S^2\log_2\left(1+\frac{v_S^2}{u_S^2}\right)\\
 &\leq v_S^2\log_2\frac1{v_S^2}+\frac{v_S^2}{\ln2}.
 \end{aligned}
\]
Thus the entropy increment in (2) is at most
$\mathcal L(d)+3q/\ln2$.  Finally use $\cU=\cH-V$; this retains the
complete upper-set asymmetry term rather than assigning signs separately
to the $B-\Gamma+\mathrm{KL}$ pieces.
\end{proof}

\section{The all-depth telescope}

Consider a finite binary restriction tree.  Every internal node $v$ queries
its own fresh uniform physical coordinate $z_v$, no such coordinate is used
at any other node, and $w_v=2^{-\operatorname{depth}(v)}$ is its reach
probability.  Terminal functions $g_\ell$ are otherwise arbitrary Boolean
functions: they may share all data coordinates and need not be tensor,
bent, orthogonal, or spectrally flat.  Let $F_v^-,F_v^+$ be the two child
subfunctions,
\[
 d_v=(F_v^--F_v^+)/2,\qquad q_v=\E d_v^2.
\]

\begin{theorem}[retained-credit tree bound]\label{thm:tree}
For every such tree, of arbitrary depth,
\begin{equation}
\begin{split}
 \cU(F_{\rm root})-\sum_{\ell}w_\ell\cU(g_\ell)
 \leq{}&\sum_\ell w_\ell V(g_\ell)-V(F_{\rm root})\\
 &+\sum_{v\ {
m internal}}w_v\mathcal L(d_v)
 +\frac3{\ln2}\,I_{\rm sel}(F_{\rm root}).
\end{split}
\tag{5}
\end{equation}
where
\[
 I_{\rm sel}(F_{\rm root})=\sum_v w_vq_v.
\]
Consequently one may drop the nonpositive root term $-V(F_{\rm root})$.
All coefficients are independent of tree depth, reach probability, terminal
bias, data-coordinate overlap, and the Fourier shapes of the terminal cores.
\end{theorem}

\begin{proof}
Apply (1) at every internal node and multiply it by $w_v$.  Since the two
child weights are $w_v/2$, both the $\cU$ terms and the $V$ terms telescope,
leaving only the root and terminal values.  Freshness of $z_v$ makes its
global physical influence exactly $w_vq_v$; summing gives the displayed
selector influence identity.
\end{proof}

This theorem covers nested rare activations: a node at depth $j$ is reached
with probability $2^{-j}$, yet neither the selector coefficient nor the
retained-packet coefficient grows with $j$.  It also permits changing,
coherent, fully overlapping terminal cores.  Each node contributes exactly
one raw difference-spectrum packet $w_v\mathcal L(d_v)$, and every internal
$V$ packet cancels exactly.  Only the probability-weighted terminal
$\cH-\cU$ packet survives.

\section{What remains unpaid}

Equation (5) is a retained-credit theorem, not a pure influence theorem.
The universal remaining question is whether
\[
 \sum_v w_v\mathcal L(d_v)
\]
admits dimension-free ownership by nonreused physical data influence and/or
terminal entropy.  The sparse bent-core surgery obstruction shows why this
packet cannot be replaced by $C\sum_vw_vq_v$: already at one node its ratio
to selector influence can diverge.  Thus the new ledger survives that
falsifier by retaining exactly the currency which the falsifier amplifies.

Repeated physical selector variables are also outside Theorem~\ref{thm:tree};
then $\sum_vw_vq_v$ need not equal actual selector influence.  Resolving
raw-packet ownership under changing cores or reused address variables is the
next unpaid step.  No claim about full FEI follows from (5).

\end{document}
